\chapter{From Power Flow to Spectral Dynamics}

This chapter derives the mathematical objects used throughout the thesis from
the electrical model upward. The purpose is pedagogical as much as technical:
the reader should see where the matrix \(L\) comes from, why the swing equation
has the form \(M\ddot\theta+D\dot\theta+L\theta=p\), and why substituting
\(e^{st}\) produces the quadratic operator \(s^2M+sD+L\). Later chapters use
these equations aggressively, so this chapter builds them slowly.

\section{The Physical Picture}

A power grid is a network of buses connected by lines and transformers. A bus
is a node where generators, loads, converters, transformers, or transmission
lines meet. At each AC bus, the voltage is not only a magnitude. It is a
sinusoid with a phase angle. In phasor notation,
\begin{equation}
  V_i=|V_i|e^{j\theta_i},
  \label{eq:guide_voltage_phasor}
\end{equation}
where \(|V_i|\) is the voltage magnitude and \(\theta_i\) is the phase angle.
Power transfer in an AC network depends strongly on differences of angles,
not on absolute angles. If every angle is shifted by the same constant, the
physics is unchanged. This is the reason a Laplacian matrix will later have
one zero eigenvalue.

The network itself is represented by the bus-admittance matrix
\begin{equation}
  Y_{\rm bus}=G+jB.
  \label{eq:guide_ybus}
\end{equation}
This matrix is built from the impedances of lines and transformers. It maps
bus voltages to injected currents:
\begin{equation}
  I=Y_{\rm bus}V.
  \label{eq:guide_current_balance}
\end{equation}
Equation~\eqref{eq:guide_current_balance} is Kirchhoff's current law in matrix
form. It is the electrical network before any dynamic-machine model is added.

\section{Complex Power and Power Flow}

The complex power injected at bus \(i\) is
\begin{equation}
  S_i=P_i+jQ_i=V_i I_i^*,
  \label{eq:guide_complex_power}
\end{equation}
where \(P_i\) is active power and \(Q_i\) is reactive power. Substituting
\(I=Y_{\rm bus}V\) gives the AC power-flow equations:
\begin{align}
  P_i(V,\theta)
  &=
  |V_i|\sum_{j=1}^n |V_j|
  \left(
  G_{ij}\cos(\theta_i-\theta_j)
  +
  B_{ij}\sin(\theta_i-\theta_j)
  \right),
  \label{eq:guide_active_power}\\
  Q_i(V,\theta)
  &=
  |V_i|\sum_{j=1}^n |V_j|
  \left(
  G_{ij}\sin(\theta_i-\theta_j)
  -
  B_{ij}\cos(\theta_i-\theta_j)
  \right).
  \label{eq:guide_reactive_power}
\end{align}
These equations determine the operating point. A power-flow solution provides
the steady voltages \(V^\star\), angles \(\theta^\star\), and injections that
balance generation and load. The dynamic model in this thesis is always a
small-signal model around such an operating point.

\begin{remark}[Where power flow enters the thesis]
The matrix \(L\) is not chosen by hand. It is the linearized active-power
response around a solved operating point. If the loading, generation dispatch,
network topology, line parameters, or converter injections change, the power
flow changes, and therefore the stiffness matrix \(L\) changes.
\end{remark}

\section{From Active-Power Jacobian to the Laplacian \texorpdfstring{\(L\)}{L}}

To obtain a dynamic model, perturb the angles around the operating point:
\begin{equation}
  \theta=\theta^\star+\Delta\theta.
  \label{eq:guide_angle_perturbation}
\end{equation}
The first-order change in active power is
\begin{equation}
  \Delta P
  \approx
  J_{P\theta}\Delta\theta,
  \qquad
  J_{P\theta}
  =
  \left.
  \frac{\partial P}{\partial\theta}
  \right|_{(V^\star,\theta^\star)}.
  \label{eq:guide_active_jacobian}
\end{equation}
In this dissertation \(P\) denotes exported electrical power in the swing
equation. With that convention, the synchronizing stiffness is this incremental
power response:
\begin{equation}
  L=J_{P\theta}
  \quad
  \hbox{on the retained angle coordinates.}
  \label{eq:guide_L_from_jacobian}
\end{equation}
If instead \(P\) is defined as net injected disturbance with the opposite sign,
the formula becomes \(L=-J_{P\theta}\). The sign convention must be fixed before
using any derivative formula.

Under the common lossless approximation \(G_{ij}\approx0\), fixed voltage
magnitudes, and sufficiently small angle differences, the off-diagonal
entries become
\begin{equation}
  L_{ij}
  =
  -|V_i^\star||V_j^\star|B_{ij}
  \cos(\theta_i^\star-\theta_j^\star),
  \qquad i\ne j,
  \label{eq:guide_L_offdiag}
\end{equation}
and the diagonal is chosen so that each row sums to zero:
\begin{equation}
  L_{ii}=-\sum_{j\ne i}L_{ij},
  \qquad
  L\ones=0.
  \label{eq:guide_L_rowsum}
\end{equation}
This is the graph Laplacian structure. The line weight is roughly
\(|V_i^\star||V_j^\star|B_{ij}\cos(\theta_i^\star-\theta_j^\star)\). A line
with high susceptance, high voltage magnitudes, and small angle separation
contributes stronger synchronizing stiffness.

\begin{example}[Two-bus stiffness]
For two buses connected by one lossless line, active power from bus 1 to bus 2
is approximately
\[
  P_{12}=|V_1||V_2|B_{12}\sin(\theta_1-\theta_2).
\]
Linearizing around \(\theta^\star\) gives
\[
  \Delta P_{12}
  \approx
  |V_1^\star||V_2^\star|B_{12}
  \cos(\theta_1^\star-\theta_2^\star)
  (\Delta\theta_1-\Delta\theta_2).
\]
Thus the line resists angle separation like a spring. That spring constant is
the line weight that appears inside \(L\).
\end{example}

\section{The Single-Machine Swing Equation}

The swing equation is Newton's second law for rotor angle. A synchronous
machine stores kinetic energy in a rotating mass. Let \(\theta_i\) be the
electrical angle of machine \(i\) relative to a synchronous reference, and let
\(\omega_i=\dot\theta_i\) be the frequency deviation. The power imbalance
accelerates the rotor:
\begin{equation}
  M_i\dot\omega_i
  =
  P_{m,i}-P_{e,i}-D_i\omega_i.
  \label{eq:guide_single_swing_speed}
\end{equation}
Here \(P_{m,i}\) is mechanical or commanded input power, \(P_{e,i}\) is
electrical air-gap power exported to the network, \(D_i\omega_i\) is a
linearized damping or primary-control term, and \(M_i\) is the inertia
coefficient in the chosen per-unit convention. Since \(\omega_i=\dot\theta_i\),
\eqref{eq:guide_single_swing_speed} becomes
\begin{equation}
  M_i\ddot\theta_i+D_i\dot\theta_i
  =
  P_{m,i}-P_{e,i}.
  \label{eq:guide_single_swing_angle}
\end{equation}

At the operating point, \(P_{m,i}^\star=P_{e,i}^\star\). The dynamic equation
therefore concerns deviations from that balance. Let \(p_i\) be a small input
disturbance and let \(\Delta P_{e,i}\) be the incremental electrical power.
Then
\begin{equation}
  M_i\Delta\ddot\theta_i+D_i\Delta\dot\theta_i+\Delta P_{e,i}=p_i.
  \label{eq:guide_single_linear_swing}
\end{equation}
The previous section showed that, to first order,
\(\Delta P_e\approx L\Delta\theta\). Therefore the electrical network supplies
the restoring term in the swing equation.

\section{The Networked Swing Model}

Stack all retained angles into \(\theta\), all inertia coefficients into
\(M=\diag(M_i)\), and all damping coefficients into \(D=\diag(D_i)\). Dropping
the \(\Delta\) notation for readability gives
\begin{equation}
  M\ddot\theta+D\dot\theta+L\theta=p.
  \label{eq:guide_network_swing}
\end{equation}
This is the model used in the fixed-\(D\) chapters. Each term has a physical
meaning:
\begin{itemize}
\item \(M\ddot\theta\): inertia resists acceleration;
\item \(D\dot\theta\): damping or primary control resists speed deviation;
\item \(L\theta\): network synchronizing stiffness resists angle separation;
\item \(p\): disturbance or control input.
\end{itemize}

Because \(L\ones=0\), the uniform angle direction must be removed by selecting
a reference bus, center-of-inertia coordinates, or a projection onto
\(\ones^\perp\). After that reduction, \(L\) is positive definite on the
retained nonzero-angle subspace if the network is connected and the operating
point is regular.

\section{Why the Exponential Ansatz Produces \texorpdfstring{\(s^2M+sD+L\)}{s2M+sD+L}}

Small-signal stability asks what happens to a tiny perturbation if no further
input is applied. Set \(p=0\) and look for modal solutions of the form
\begin{equation}
  \theta(t)=ve^{st},
  \label{eq:guide_modal_ansatz}
\end{equation}
where \(v\) is the mode shape and \(s\) is the complex growth rate. Then
\begin{equation}
  \dot\theta(t)=sve^{st},\qquad
  \ddot\theta(t)=s^2ve^{st}.
  \label{eq:guide_exponential_derivatives}
\end{equation}
Substituting \eqref{eq:guide_modal_ansatz} into
\eqref{eq:guide_network_swing} gives
\[
  Ms^2ve^{st}+Dsve^{st}+Lve^{st}=0.
\]
Since \(e^{st}\ne0\), it can be divided out:
\begin{equation}
  (s^2M+sD+L)v=0.
  \label{eq:guide_qep_from_swing}
\end{equation}
This is the quadratic eigenvalue problem (QEP). It is not an arbitrary
mathematical object. The factor \(s^2\) comes from acceleration, \(s\) comes
from velocity damping, and the constant term comes from the network restoring
force.

\begin{table}[ht]
\centering
\caption{Physical origin of the QEP terms.}
\label{tab:qep_origin}
\begin{tabular}{lll}
\toprule
\rowcolor{softblue}
\thead{Term} & \thead{Physical source} & \thead{Engineering interpretation}\\
\midrule
\(s^2M\) & rotor or virtual mass & acceleration and inertia\\
\(sD\) & damping, droop, primary response & decay of speed deviations\\
\(L\) & linearized active-power network & synchronizing stiffness\\
\((sI-A_{cc})^{-1}\) & eliminated controller states & dynamic memory of controls\\
\bottomrule
\end{tabular}
\end{table}

\section{How to Read a Complex Pole}

If a root is \(s=\sigma+j\omega\), the perturbation contains
\[
  e^{st}=e^{\sigma t}e^{j\omega t}.
\]
The real part \(\sigma\) controls decay or growth. The imaginary part
\(\omega\) controls oscillation frequency. Stability requires
\(\sigma<0\). The damping ratio is
\begin{equation}
  \zeta=-\frac{\sigma}{\sqrt{\sigma^2+\omega^2}}.
  \label{eq:guide_zeta}
\end{equation}
This number measures how quickly the oscillation decays relative to its
frequency. A line upgrade can increase \(\omega\) by stiffening the network,
but if \(\sigma\) does not become more negative, \(\zeta\) may decrease. This
is the basic reason graph stiffness and damping margin must be kept separate.

\section{Where Inverters Enter the Same Chain}

An inverter-based resource can enter this chain at several levels. In a simple
reduced model, a grid-forming inverter may be represented by effective
inertia, damping, or droop terms added to \(M\) and \(D\). A grid-following
inverter may alter the operating point and the Jacobian through its injection
and controls. In a detailed small-signal model, inverter controllers add
states: PLL states, current-loop states, voltage-loop states, filters, limits,
and plant controls. If those states are retained, they appear in the full
state matrix. If they are mathematically eliminated, they appear through a
Schur complement and produce the rational self-energy \(\Pi(s)\) developed in
the self-energy part of the thesis. The dynamic damping object \(\Sigma(s)\)
is the special case where
the rational correction enters through the velocity channel.

Thus the thesis path is:
\begin{equation}
  \hbox{power flow}
  \rightarrow
  L
  \rightarrow
  M\ddot\theta+D\dot\theta+L\theta=p
  \rightarrow
  s^2M+sD+L
  \rightarrow
  s^2M+sD_0+L+\Pi(s).
  \label{eq:guide_thesis_path}
\end{equation}
Every later result is a refinement of one arrow in this chain.

\begin{figure}[ht]
\centering
\begin{tikzpicture}[
  stage/.style={
    rectangle,
    rounded corners=3pt,
    draw=#1!70!black,
    fill=#1!10,
    text width=0.25\textwidth,
    minimum height=1.15cm,
    align=center,
    font=\small
  },
  flow/.style={-{Latex[length=2.2mm]}, thick, draw=thesisgray},
  note/.style={font=\scriptsize\color{thesisgray}, align=center}
]
\node[stage=thesisblue] (pf) {\textbf{Power-flow solution}\\
\(V^\star,\theta^\star,Y_{\rm bus}\)};
\node[stage=thesisgreen, right=0.65cm of pf] (jac) {\textbf{Active-power Jacobian}\\
\(L=J_{P\theta}\)};
\node[stage=thesisorange, right=0.65cm of jac] (swing) {\textbf{Swing layer}\\
\(M\ddot\theta+D\dot\theta+L\theta=p\)};
\node[stage=thesisblue, below=0.85cm of jac] (qep) {\textbf{Fixed-\(D\) spectral problem}\\
\((s^2M+sD+L)v=0\)};
\node[stage=thesisred, right=0.65cm of qep] (nep) {\textbf{Condensed-controller problem}\\
\(T(s)=s^2M+sD_0+L+\Pi(s)\)};
\draw[flow] (pf) -- (jac);
\draw[flow] (jac) -- (swing);
\draw[flow] (swing.south) -- ++(0,-0.42) -| ([xshift=0.32cm]qep.north);
\draw[flow] (jac.south) -- (qep.north);
\draw[flow] (qep) -- (nep);
\node[note, below=0.18cm of pf] {electrical constraints};
\node[note, below=0.18cm of qep] {modal damping screen};
\node[note, below=0.18cm of nep] {controller memory};
\end{tikzpicture}
\caption{Modeling chain from AC power flow to the rational NEP. The same
network operating point that sets voltages and flows also defines the
stiffness operator \(L\). Controller condensation changes the operator class,
because the damping contribution becomes frequency dependent.}
\label{fig:prelim_modeling_chain}
\end{figure}

\FloatBarrier

\section{Linearization and Spectral Parameter}

Let a power-system dynamic model be written locally as a differential-algebraic
system,
\begin{align}
  \dot x &= f(x,y,\rho),\\
  0 &= g(x,y,\rho),
\end{align}
where \(x\) are dynamic states, \(y\) are algebraic variables, and \(\rho\) is
a parameter or action. At a regular equilibrium
\((x^\star,y^\star)\), algebraic variables can be eliminated locally when
\(g_y\) is nonsingular, giving a reduced state matrix
\begin{equation}
  A_{\rm red}=f_x-f_y g_y^{-1}g_x.
  \label{eq:dae_reduced_A}
\end{equation}
Small-signal stability is determined by the eigenvalues of \(A_{\rm red}\).
The spectral parameter is denoted by \(s\in\C\), and a pole is stable when
\(\Real(s)<0\).

For an oscillatory pole \(s=\sigma+j\omega\), the frequency and damping ratio
are
\begin{equation}
  f=\frac{|\omega|}{2\pi},\qquad
  \zeta=-\frac{\sigma}{\sqrt{\sigma^2+\omega^2}}.
  \label{eq:prelim_zeta}
\end{equation}
The damping-ratio derivative satisfies
\begin{equation}
  d\zeta
  =
  -\frac{\omega^2}{(\sigma^2+\omega^2)^{3/2}}\,d\sigma
  +
  \frac{\sigma\omega}{(\sigma^2+\omega^2)^{3/2}}\,d\omega,
  \label{eq:prelim_dzeta}
\end{equation}
so a frequency shift alone is not a damping certificate.

\section{Graph Laplacians and Synchronizing Stiffness}

Under a lossless active-power approximation, the synchronizing stiffness matrix
has Laplacian structure. For \(i\ne j\),
\begin{equation}
  L_{ij}=-V_iV_j b_{ij}\cos(\theta_i^\star-\theta_j^\star),
  \label{eq:prelim_lap_offdiag}
\end{equation}
with diagonal entries chosen so that
\begin{equation}
  L\ones=0.
\end{equation}
If the network is reciprocal and the operating point is sufficiently regular,
\(L=L^T\succeq0\). The zero eigenvalue corresponds to uniform angle shift and
is removed by fixing a reference or projecting onto \(\ones^\perp\).

Let \(M=\diag(M_i)\succ0\). The inertia-normalized stiffness is
\begin{equation}
  \Lt=M^{-1/2}LM^{-1/2}.
  \label{eq:prelim_ltilde}
\end{equation}
On the nonzero subspace, write
\begin{equation}
  \Lt Q=Q\Lambda,\qquad Q^TQ=I,\qquad
  \Lambda=\diag(\nu_1,\ldots,\nu_n),\quad \nu_k>0.
  \label{eq:prelim_graph_modes}
\end{equation}
The columns of \(Q\) are graph-inertia modes. They diagonalize stiffness. They
do not automatically diagonalize damping.

\section{Second-Order Systems and QEPs}

The retained second-order model is
\begin{equation}
  M\ddot\theta+D\dot\theta+L\theta=p.
  \label{eq:prelim_second_order}
\end{equation}
Substituting \(\theta(t)=ve^{st}\) gives
\begin{equation}
  P(s)v=(s^2M+sD+L)v=0.
  \label{eq:prelim_qep}
\end{equation}
This is a QEP. The corresponding first-order linearization is
\begin{equation}
  \mathcal A=
  \begin{bmatrix}
    0&I\\
    -M^{-1}L&-M^{-1}D
  \end{bmatrix},
  \label{eq:prelim_first_order}
\end{equation}
but the second-order form should be retained conceptually because the matrices
\(M\), \(D\), and \(L\) have distinct physical meanings.

In inertia-normalized variables \(\vartheta=M^{1/2}\theta\),
\begin{equation}
  \ddot\vartheta+\Dt\dot\vartheta+\Lt\vartheta=M^{-1/2}p,\qquad
  \Dt=M^{-1/2}DM^{-1/2}.
  \label{eq:prelim_normalized}
\end{equation}
The modal damping matrix is
\begin{equation}
  \Gamma=Q^T\Dt Q.
  \label{eq:prelim_gamma}
\end{equation}
The equality \(Q^T\Lt Q=\Lambda\) is automatic by construction; the equality
\(Q^T\Dt Q=\diag(\cdot)\) is an additional damping assumption.

\section{Simple Eigenvalue Sensitivity}

For a polynomial or nonlinear eigenvalue problem
\begin{equation}
  T(s,\rho)v=0,\qquad u^HT(s,\rho)=0,
\end{equation}
with a simple eigenvalue \(s_\star\), the first derivative is
\begin{equation}
  \frac{ds_\star}{d\rho}
  =
  -\frac{u^H T_\rho(s_\star,\rho)v}
  {u^H T_s(s_\star,\rho)v}.
  \label{eq:prelim_nep_sens}
\end{equation}
For a QEP \(T(s)=s^2M+sD+L\), this becomes
\begin{equation}
  \frac{ds_\star}{d\rho}
  =
  -\frac{u^H(s_\star D_\rho+L_\rho)v}
  {u^H(2s_\star M+D)v}
  \label{eq:prelim_qep_sens}
\end{equation}
when \(M\) is fixed. Equation~\eqref{eq:prelim_qep_sens} is central because it
shows that line actions and damping actions enter the numerator differently,
while the denominator depends on the mode and the damping operator.

\section{Schur Complements}

Partition a full first-order linearization into retained states \(x_e\) and
condensed states \(x_c\):
\begin{equation}
  A=
  \begin{bmatrix}
    A_{ee} & A_{ec}\\
    A_{ce} & A_{cc}
  \end{bmatrix}.
  \label{eq:prelim_partitioned_A}
\end{equation}
The eigenvalue equation \((sI-A)x=0\) is equivalent, when \(sI-A_{cc}\) is
invertible, to
\begin{equation}
  \left[sI-A_{ee}-A_{ec}(sI-A_{cc})^{-1}A_{ce}\right]x_e=0.
  \label{eq:prelim_schur}
\end{equation}
This is an exact nonlinear eigenvalue problem for the retained states. It does
not approximate the eliminated states; it represents them through their
resolvent.

\section{Rational Matrix Functions and Self-Energy}

A rational matrix function has the form
\begin{equation}
  R(s)=R_0+\sum_{\ell=1}^m \frac{R_\ell}{s-\alpha_\ell}
  \label{eq:prelim_rational}
\end{equation}
after partial-fraction expansion when the poles are simple. The residues
\(R_\ell\) determine how strongly each eliminated pole \(\alpha_\ell\) couples
to the retained coordinates. In the Schur operator, those residues are
generated by \(A_{ec}\), \(A_{ce}\), and the eigenvectors of \(A_{cc}\).

In a second-order retained representation, the rational contribution appears as
a self-energy,
\begin{equation}
  T(s)=s^2M+sD_0+L+\Pi(s),
  \qquad
  \Pi(s)=G(sI-A_{cc})^{-1}(B_\theta+sB_\omega).
  \label{eq:prelim_sigma}
\end{equation}
The term self-energy is useful because it emphasizes that the retained model is
not isolated. It carries the dynamic influence of states that have been
mathematically eliminated. When \(\Pi(s)=s\Pi_D(s)\), the dynamic damping
matrix is \(\Sigma(s)=D_0+\Pi_D(s)\).

\section{Certificates, Screens, and Gates}

The thesis uses three words with specific meanings.

\begin{definition}[Certificate]
A certificate is a mathematical or numerical condition that directly implies a
claimed margin property under stated assumptions.
\end{definition}

\begin{definition}[Screen]
A screen is a lower-cost indicator that may identify candidates for action or
inspection but does not imply the claimed margin property by itself.
\end{definition}

\begin{definition}[Gate]
A gate is an accept/reject test that a proposed certificate or screen must pass
before it is used for a stronger claim.
\end{definition}

The distinction matters throughout the dissertation. A graph eigenvalue can be
a certificate in a proportional or commuting fixed-\(D\) theory, a screen in a
non-commuting fixed-\(D\) theory, and an incomplete descriptor after controller
condensation. The same number changes epistemic status when the operator class
changes.
